\noindent
{\Large \textbf{Titre :} « \textbf{RECONSTRUCTION D'IMAGES ISSUES DE TOMODENSITOMÉTRIE} » }\vspace{0.5cm}\\
\textbf{Auteur :} ANDRIANARIVONY Selison Fréderic \\
\textbf{Nombre de pages :} 120 \\
\textbf{Nombre de figures :} 22 \\
\textbf{Nombre de tableaux :} 2
\vspace{0.5cm}

\noindent
\paragraph{RESUMÉ} \\
Ce mémoire de master à visée de recherche s'intéresse à la reconstruction d'images issues de la tomodensitométrie, une technique d'imagerie médicale essentielle pour l'observation non invasive de l'intérieur du corps humain. L'objectif est d'améliorer la qualité des images obtenues à partir des données brutes issues des scanners, en explorant des méthodes avancées de reconstruction qui permettent de réduire le bruit, d'augmenter la résolution et de limiter la dose de rayonnement administrée aux patients.\\

Le travail s'inscrit dans une démarche scientifique visant à proposer des approches innovantes, probablement basées sur des modèles mathématiques, des algorithmes d'optimisation ou des techniques d'apprentissage automatique. En évaluant les performances de ces méthodes sur des données réelles ou simulées, on cherche à contribuer à l'amélioration des pratiques cliniques et à l'évolution des systèmes d'imagerie médicale vers plus de précision et de sécurité pour les patients.
\vspace{0.5cm}

\noindent
\textbf{ENCADREUR DE MEMOIRE :} Madame RAMAFIARISONA Malalatiana

\vspace{0.5cm}

\noindent
\textbf{ADRESSE :} Bloc 54 CU Ambohipo \\
\textbf{CONTACTS :} +261 34 77 161 84 \\
\textbf{Email :} frederic.andrianarivo@gmail.com
